\itemhead{习题4.2 · 第3题（教材p.113）}{XJ-B1-C04-HW-P113-Y1144}
\label{XJ-B1-C04-HW-P113-Y1144}
{\footnotesize 湘教版教材p.113 / PDF第120页 / 习题4.2 / 第3题}
\field{原题}{一辆新车的售价为16万元，使用 $x$ 年后，其价值 $P$（万元）为 $P=16\cdot0.9^x$。
\subq{1}说明式中0.9的含义；
\subq{2}求这辆车使用3年、9年后的价值。}
\field{参考答案}{（1）每年的保留价值为前一年的90\%，年折旧率10\%；（2）$P(3)=11.664$ 万元，$P(9)=6.198727824\approx6.20$ 万元。}
\field{核对说明}{近似值差异：16×0.9\textasciicircum{}9=6.198727824万元，通常保留两位小数为6.20；答案册PDF32写约6.19，现采用精确值舍入。}
{\small \field{关联知识}{E03 指数增长、衰减与实际模型（教材帮p.131、132、154 / PDF140、141、163）}
\field{方法与易错点}{确认初始量、时间单位、倍率与时刻起点，再代入或逐次比较阈值。 易错：把年折旧10\%写成乘0.1；“超过”用成“达到”；时间起点错一位；把衰减至接近0当实际为0。}
\field{答案依据}{答案册相关题组：PDF第32页（册内页31）；按当前题库原题逐项重算，不按旧版题号机械对套。}}
